\documentclass{article}
\usepackage[utf8]{inputenc}
\usepackage{parskip}
\usepackage{amsmath}
\usepackage{amssymb}
\usepackage[thinc]{esdiff}

\title{Solving Linear Second Order ODEs}

\author{
    Yip, Henry\\
    s2231321@ed.ac.uk
    \and
    Richmond, Lorian\\
    s2188785@ed.ac.uk
}

\date{March 2022}

\begin{document}

\maketitle
\begin{abstract}
    This paper includes
    \begin{itemize}
     \item Several proofs behind choosing the Ansatz: $y=e^{\lambda x}$
     \item Generalizing solutions for Linear Second Order ODEs using Linear Algebra
     \item Physical applications in Quantum Mechanics and Classical Mechanics
     \end{itemize}
\end{abstract}
\section{Introduction}

Consider the homogeneous equation:
\begin{equation}
\ddot{y}+ \alpha\dot{y}+ \beta y = 0
\end{equation}
By basic differentiation techniques, we have 
 \begin{align}
 \frac{d}{dx}(e^x)=e^x\\
 \implies \frac{d}{dx}(ae^x)=ae^x
 \end{align}
 
 Suppose that f: $\mathbb{R} \rightarrow \mathbb{R}$ is differentiable everywhere and $f'(x)=af(x)$.
 There are no other functions in which the derivative is equal to itself. 
 \subsection{Proof 1}

 If we let $g(x)=\frac{af(x)}{e^{ax}}$


$ f(x)$ and $e^{ax}$ are differentiable $\rightarrow g'(x)$ is defined.
\begin{align}
     g'(x)&=\frac{f'(x)e^{ax}-af(x)e^{ax}}{(e^{ax})^2}\\
     g'(x)&=\frac{af(x)e^{ax}-af(x)e^{ax}}{(e^{ax})^2}\\
     g'(x)&=0
\end{align}
 By identity theorem, $g'(x)=0\rightarrow g(x)= c$, where $c\in \mathbb{R}$\\
 \implies $f(x)=Ae^{ax}$, where $A=ac$. We take $a=1$ and hence $A=c$.
 The above statement can be proved by Picard–Lindelöf theorem, however we won't cover it here.
 \subsection{Proof 2}
 Suppose that f: $\mathbb{R} \rightarrow \mathbb{R}$ is differentiable everywhere and $f(x)=f'(x)$
 \begin{align}
     1&=\frac{f(x)}{f'(x)}\\
     x+c&=\int{\frac{f(x)}{f'(x)}}, c\in \mathbb{R}\\
     x+c&=ln(f(x))\\
     f(x)&=e^{x+c}\\
     f(x)&=\gamma e^x, \gamma =e^c
 \end{align}
 \subsection{Ansatz}
 Therefore, we take the Ansatz: $y=e^{\lambda x}$ and obtain:
 \begin{align}
  \ddot{y}+ \alpha\dot{y}+ \beta y &= 0\\
     ({\lambda}^2 + \alpha\lambda  + \beta)(e^{\lambda x}) &=0
\end{align}
     Obviously $e^{\lambda x}=0$ is undefined, so 
     \begin{equation}
     {\lambda}^2 + \alpha\lambda  + \beta=0
     \end{equation}
     \begin{equation}
     \lambda_1=\frac{a}{2}+\sqrt{\frac{a^2}{4}-b} 
     \end{equation}
     \begin{equation}
     \lambda_2=\frac{a}{2}-\sqrt{\frac{a^2}{4}-b}
     \end{equation}
     
     We can imagine both $\lambda_1$ and $\lambda_2$ satisfies the equation, so $y$ is a linear combination of $e^{\lambda_1 x}$ and $e^{\lambda_2 x}$, we can rewrite this as:
     \begin{equation}
         y=Ae^{\lambda_1 x}+Be^{\lambda_1 x}
     \end{equation}
     Where A \& B $\in \mathbb{R}$
     
     \bf{A Rigorous Statement:}
     Let S be the solution space of 
     $$\ddot{y}+ \alpha\dot{y}+ \beta y &= 0$$
     If $\lambda_1 \ne \lambda_2$, then $\{e^{\lambda_1 x} , e^{\lambda_2 x}\}$ is a basis for S.
     
     Note that the above is only true for $\lambda_1 \ne \lambda_2$. For $\lambda_1 = \lambda_2$,$\{e^{\lambda_1 x} , te^{\lambda_1 x}\}$ is a basis for S. The proof will be written some other time.
 \end{document}

